% !TeX program = xelatex
% ===========================================================================
%  Dislocation and Reflection on the ``Non-Quantifiable'' Work of Music
%   -- English version; layout supplied uniformly by musepaper-en.sty
%
%  Compile: xelatex paper.en.tex   (three passes for cross-references)
% ===========================================================================
\documentclass[11pt,a4paper]{article}

\usepackage{musepaper-en}

% Running head
\runninghead{Dislocation and Reflection on the ``Non-Quantifiable'' Work of Music}

\title{Dislocation and Reflection on the ``Non-Quantifiable'' Work of Music}
\author{WANG Zijian}
\date{}

%% ---------------------------------------------------------------------------
%%  Citation practice (per the documentation standard of the Journal of the
%%  Central Conservatory of Music): all cited literature appears in page
%%  footnotes, with no reference list appended at the end of the paper.  The
%%  original \begin{thebibliography} has therefore been replaced by \footnote
%%  at each point of citation.
%%  NOTE: the standard requires the cited page number; wherever 【待核】 /
%%  【X】 appeared in the Chinese source it is preserved verbatim below as
%%  [to be verified] / [X], since no edition was available to check it against.
%% ---------------------------------------------------------------------------
\begin{document}

% The title page carries neither running head nor page number (standard
% practice for journal submissions); page numbers begin on page 2.

% ===========================================================================
% Title page
% ===========================================================================
\begin{center}
  {\sffamily\bfseries\fontsize{17}{21}\selectfont
   Dislocation and Reflection on the ``Non-Quantifiable'' Work of Music\par}
  \vspace{0.9em}
  {\normalsize WANG Zijian\par}
  \vspace{0.35em}
  {\small (MozunStudio, Hefei 230000, Anhui, China)\par}
  \vspace{0.25em}
  {\small Correspondence: 368164331@qq.com\par}
\end{center}
\vspace{0.4em}

\begin{abstracten}
The proposition that ``music is non-quantifiable'' has long been treated as a default
premise of the music profession, and music has been habitually assigned to the categories
of ``the art of humanity'' and ``the perception of the soul''. This paper argues that this
proposition should first be made precise: is it ``music'' that cannot be quantified, or is
it ``musical emotion'' that cannot be quantified? Having defined quantization as ``lossy
compression through discretization and symbolization'', this paper proposes four criteria
for quantifiable objects, of which the ternary relation ``object $\times$ purpose $\times$
tolerance'' constitutes the core. On this basis it argues that music comprises two
layers---the static and the dynamic: the layer of form and structure is static and
quantifiable, whereas the layer of aesthetic appreciation and human-made relative dynamics
is non-quantifiable, the two coexisting in a collaborative mode exemplified by
``accompaniment track and human voice''. This paper further re-examines the definition of
music, proposing that the ``regularity'' of music is ``the proportion of sound occupying a
given period of time'', thereby establishing the ontological status of the temporal
dimension in music. Citing Ji Kang's \textit{Sheng Wu Ai Yue Lun} (On the Absence of
Sorrow and Joy in Sound) and Hanslick's theory of autonomy, it corroborates the
human-made attribute of the aesthetic layer; and, with Zhu Zaiyu's twelve-tone equal
temperament, Bach's \textit{The Well-Tempered Clavier}, the evolution of music notation,
and the mechanism of composition (the isomorphism of motif and training), it argues for
the quantifiability of composition and arrangement. Finally, this paper reveals the
institutional roots of the absence of quantification consciousness---the model of oral
transmission, the dislocation of ``right/wrong'' evaluation, and the imbalance of
disciplinary construction---and looks forward to the significance for AI music arrangement
of a layered quantization scheme (absolute coordinates at the structural layer, explicit
tolerance at the aesthetic layer).
\end{abstracten}

\keywords{music quantification; ontology of sound; temporal structure;
aesthetics of music; layered quantization; artificial intelligence composition}

\newpage
% ===========================================================================
% Table of contents
% ===========================================================================
{\sffamily\bfseries\fontsize{14}{17}\selectfont
 \begin{center}Contents\end{center}}
\vspace{0.6em}
\renewcommand{\contentsname}{}
\vspace{-3.2em}
\tableofcontents
\newpage

% ===========================================================================
% Body
% ===========================================================================
\section{Introduction: Statement of the Problem}

In the traditional conception of music, ``music is non-quantifiable'' is an almost
self-evident axiom. We habitually assign music to ``the art of humanity'' and ``the
perception of the soul'', as though music would cease to be music the moment a scale and a
grid were applied to it. This preconception is so firmly entrenched that the entire
profession's digitalization---from notation software to the digital audio workstation
(DAW)---has had to proceed cautiously, under the defense of ``artistry''.

Yet the present author, for all his shallowness, holds that this axiom itself deserves to
be questioned: when we say that ``music is non-quantifiable'', is it ``music'' that cannot
be quantified, or is it ``musical emotion'' that cannot be quantified? These two
propositions point to entirely different objects: the former concerns the ontological
attributes of the organization of sound, the latter the psychological effects upon the
listening subject. To conflate the two is precisely the most fundamental dislocation in the
``non-quantifiable'' work of music.

The argumentative path of this paper is as follows: first, it defines ``quantization'' and
the criteria for quantifiable objects; it then argues that music comprises a quantifiable
and a non-quantifiable layer, and that the two coexist collaboratively; it goes on to
re-examine the definition of music and establish the ontological status of the temporal
dimension; thereafter, taking music history and compositional practice as evidence, it
examines the quantifiability of composition and arrangement; and finally it reveals the
institutional roots of the absence of quantification consciousness and looks forward to a
feasible scheme of layered quantization.

\section{What Is Quantization: Conceptual Definition and Criteria}

\subsection{The Definition and Structure of Quantization}

From a professional standpoint, quantization $=$ the process of mapping a continuous (or
high-cardinality) set of objects onto a discrete (countable) set of symbols. It is usually
composed of three structural components. First, the source set---a continuous domain or a
high-resolution domain, such as a continuous stretch of time or a continuous axis of
pitch. Second, the partition---cutting the source set into mutually disjoint units, each
unit being called a ``bin''. Third, the mapping rule---each bin is assigned a
representative symbol (usually its center or centroid), and thereafter all elements
falling into that bin are referred to by that symbol.

Quantization is therefore never one thing but a conspiracy of two: discretization
(cutting) plus symbolization (naming). The cutting must also yield something that can be
called by name before quantization is complete. And its essence is a kind of ``lossy
compression''---compressing the uncountable into the countable, at the price of flattening
out the differences within each bin. Indeed, human cognition is itself a grand
quantization: the senses are a continuous stream, concepts are the bins that have been
cut, and the word ``red'' is nothing more than the centroid symbol of a million spectra of
red.

\subsection{Criteria for Quantifiable Objects}

Anything quantifiable must simultaneously satisfy four criteria. First, there must exist a
structure of difference---there must be comparable distances or relations of equivalence
among the objects; where not even ``resemblance'' can be spoken of, there is nothing to
cut. Second, there must exist topology/adjacency---the partition must respect
``neighborhood''; things within a bin should be more similar to one another than to things
in other bins, otherwise the partition is arbitrary hacking. Third, there must be a
tolerable tolerance---the differences within a bin are immaterial to the given purpose.
This criterion is the most crucial: quantifiability is never an attribute of the object
itself, but a relation among ``object $\times$ purpose $\times$ tolerance''. Fourth, the
partition must be stable and reproducible---quantizing the same object ten times must
yield the same symbol, and the mapping rule must be a determinate function.

Strictly speaking, then, there is no such thing as ``a naturally quantifiable thing'';
there is only ``an object that, under a certain tolerance and for a certain purpose, can be
finitely partitioned''.

\subsection{The Pathology of the Non-Quantifiable}

Turning to the non-quantifiable, its pathology is equally clear and has roughly three
aspects. First, difference is equally fatal at every scale---no tolerance can be safely
applied. Zeno's arrow paradox is its type case: cut motion into instants and motion
vanishes on the spot; the very ``aliveness'' of motion lies in continuity itself. Second,
meaning lives in the whole, not in the part---to cut is to kill, as with the breathing of
a musical phrase or the savor of a line of poetry. Third, there is no stable
identity---each occurrence is different, so it can neither be compared nor classified.

\subsection{Four Re-Positionings of the Non-Quantifiability Thesis}

The three aspects listed in \S2.3 answer the question ``what is non-quantifiable?''; what
is to be dealt with here is another question: how far each of the several common arguments
that assert ``music is non-quantifiable'' ``holds, and up to which layer''. The two are
two sides of the same coin, yet their conclusions differ---the former speaks of the
properties of the object, the latter of the boundary of an argument. Once each is
re-positioned, the layered thesis of this paper acquires its entry point.

\textbf{First, the continuum argument.} This view holds that quantization presupposes
partition, whereas the things that matter most in music---portamento, crescendo, rubato
---are all continua; once they are cut into discrete units, continuity itself is lost. The
part of this argument that holds is incontestable: a continuum cannot be discretized
``losslessly''. But it then oversteps its boundary---``something will be lost'' does not
equal ``it cannot be cut''; it only means that one must ask further: what exactly is lost,
and does it fall within the tolerance of the given purpose? This question is precisely the
substance of criterion three. Indeed, human cognition is itself a grand quantization: what
the senses receive is a continuous spectrum, and the word ``red'' is merely one symbol
into which a million kinds of red have been compressed. Were the continuum argument
tenable, the word ``red'' ought not to exist either.

\textbf{Second, the holistic argument.} This view holds that the meaning of music lives in
an entire phrase or even an entire work, and that to cut is to kill, as with the breathing
of a phrase or the savor of a melody. This argument conflates ``the attribution of
meaning'' with ``the existence of the object''. A musical phrase can perfectly well be
recorded in its entirety as a set of coordinates with definite starting and ending
instants---each of its notes has a determinate position, and this does not prevent it from
being grasped by the listener as a whole; the savor is an effect produced in the listener
when confronted with this set of coordinates, and it belongs to the receiving end, not to
the organization of sound itself. Wholeness does not exclude coordinatizability; it merely
imposes one demand upon the coordinate system: it must be able to accommodate
``relations'', and not merely points that are isolated from one another.

\textbf{Third, the subjectivity argument.} This view holds that the emotions aroused by
music are private and incommensurable, and hence cannot be quantified. This argument is
precisely the specimen of the dislocation pointed out at the opening of this paper---it
discusses ``the psychology of the listener'' but uses that to pass judgment on ``music''.
Whether emotion is quantifiable is another question (psychology measures it along two
dimensions, valence and arousal, which is itself an attempt at quantification), but
whatever the answer, it cannot be used to infer backwards that the organization of sound is
unmeasurable. To put it more cautiously: even if one concedes that the world of emotion
refuses a scale, that only draws a boundary for the aesthetic layer and has nothing to do
with the structural layer.

\textbf{Fourth, the historicity argument.} This view holds that the value of music shifts
with the times, that yesterday's dissonance is today's norm, and that consequently no
stable quantity exists. This argument refutes ``the constancy of aesthetic standards'', not
``the measurability of structure''. The same major triad in 1722 and today: its pitch
relations, its ratios of duration, and its spectral structure are no different; what
changes is the ``evaluation'' of it, not the thing itself. Historicity is not only no
obstacle, it precisely confirms the layered thesis of this paper: the structural layer is
stable, the aesthetic layer flows.

Viewed together, the four arguments share one and the same structure: each holds at ``a
certain layer'', yet each is used to support a ``cross-layer'' conclusion. The reason this
transgression long went unnoticed is that the word ``music'' was never taken apart in the
discussion---it refers simultaneously to the organization of sound and to auditory
experience, so that a claim about the latter was smoothly read as a claim about the former.
Once these four arguments are each re-positioned in the aesthetic layer, they no longer
conflict with the quantifiability of the structural layer; the layering proposed in
Section 3 of this paper exists precisely to give each of them the place where it belongs.

\section{The Dialectic of Musical Quantifiability: Stasis and Motion}

\subsection{A Rebuttal of the ``Performance Is Unrepeatable'' Thesis}

Following the logic above, someone may ask directly: music is indisputably a
non-quantifiable thing, so what are you discussing? It is true that even the finest
performer cannot produce two identical performances---this is physically impossible. But
does that prove that music is non-quantifiable? The present author thinks not. If music
really were non-quantifiable, how then is one to explain the efficacy of a DAW and of any
notation software in recognizing and locating notes? This is an obvious fact, and it cannot
be dismissed with a blanket judgment. In philosophical terms, we must see both the static
side of things and the dynamic side of things.

\subsection{The Layering of Music: Structural Layer and Aesthetic Layer}

What, then, is stasis for music, and what is motion? The present author's conclusion is
this: with respect to ``the form and structure of music'', music is static and
quantifiable; with respect to ``the aesthetics of music'' and to ``human-made relative
dynamics'', music is non-quantifiable.

This is not difficult to understand. Suppose you record a note in a DAW or in notation
software; the next time it is played back, that note will be sitting exactly where it
belongs, obediently; whereas if a human plays it, a slight deviation may appear---and it is
precisely this ``deviation'' that embodies ``humanization''. The drift between the key and
the grid is exactly the trace of the performer's breathing.

\subsection{The Collaboration of Stasis and Motion: The Case of ``Accompaniment Track and Human Voice''}

Must the two, then, be irreconcilable? The answer is obviously no. Anyone who has sung
knows the experience of singing along with an accompaniment track: since the track is
static, the singer spontaneously moves relative to it in order to maintain a relative
balance; and between a live piano accompanist and a singer there likewise exists a tacit
balance, failing which the whole performance would simply fall apart. It follows that
stasis is not the enemy of motion but the coordinate system in which motion can occur and
remain ordered: it is precisely because the grid is constant that deviation becomes
meaningful; it is precisely because coordinates exist that breathing has a place. Stasis
and motion each have their own office and each completes the other.

\section{A Re-Examination of the Ontology of Music: The Temporal Dimension}

\subsection{A Review of the General Definition of Music}

The descriptions of music in music-theory textbooks are not wrong, but for a long time we
have overlooked one key question: what is music? General explanations run, for example,
that ``music is an art form and cultural activity that expresses thought, emotion and
social life through the organization of regular sound waves'', or that ``music is an
auditory art constituted from the elements of pitch, dynamics, duration and timbre through
the forms of rhythm, melody and harmony''. These accounts are all correct, but note one
phrase in them---``the organization of regular sound waves'': what, then, is this
``regularity''? And as for the so-called ``use of the elements of pitch, dynamics, duration
and timbre'', what is this series of ``arts''?

\subsection{Regularity: The Proportion of Sound Occupying a Given Period of Time}

The present author's view is this: this regularity is ``the proportion of sound occupying a
given period of time''. The sentence looks abstract but is in fact easy to grasp: the
duration of a concert is always relatively fixed---two hours' worth of content, even with
small errors, cannot be compressed into half an hour unless an abnormal accident occurs.
And the pitches, the rests, and the different instruments within it are nothing more than
``quantifiable subsets within time''.

From this one may draw a conclusion that has long been overlooked: time is the unit that
truly possesses ontological status in music. All sonic elements---pitch (the distribution
of frequency within a unit of time), timbre (the spectrum evolving over time), harmony (the
vertical slice at a single instant), rhythm (the distribution of density of sound along
the time axis)---can all be projected onto the time axis. Music is not a set of sounds, but
the organization of sounds in time.

\subsection{The Distinction between the Ontology of Music and the ``Human'' Carrier}

Someone may say that such a definition would break the ``aesthetics'' and ``artistry'' of
music. The present author holds that this worry does not stand, that the two do not
conflict, and that this is a simple question long overlooked. Let us shift our angle: the
existing sound-source systems and DAW arranging equipment are already quite
sophisticated---so do the musical works an arranger produces with a DAW really lack
``aesthetics'' and ``artistry''? Obviously not. Since the DAW itself is a tool for
quantifying music that is acknowledged as such by the profession, whence comes the talk of
``breaking aesthetics and artistry''? We cannot say that listening to music through
headphones does not count as appreciation while attending a concert does---that is plainly
illogical, because the ``ontology of music'' remains the medium of ``sound''; whether it is
produced by a human concerns only ``the process itself'', not ``the final result''.

The most blunt way to put the question is this: when the audience pays, are they
``appreciating the music'' or ``appreciating the producer of the music''? If music were
defined solely by its producer, then indeed this misconception could be explained---the
so-called lack of ``aesthetics'' and ``artistry'' is nothing but the result of aesthetics
having been bound to the carrier ``human''.

Anyone who has studied philosophy will be familiar with this sentence: regularity does not
shift with the human carrier; it can be actively made use of, but it by no means shifts
with individual will. By the same token, since music has long been defined as ``the
organization of regular sound waves'', does imposing the carrier ``human'' upon that
regularity not appear excessively forceful, and contrary to philosophy itself?

\section{On the Absence of Sorrow and Joy in Sound: The Human-Made Attribute of the Aesthetic Layer}

\subsection{The Anticipation in Ji Kang's \textit{Sheng Wu Ai Yue Lun}}

If one starts from ``the form of quantization'', then ``music is the manifestation of the
soul''---at first hearing there is nothing wrong with this, but it overlooks one point:
the human psyche has nothing to do with music itself. This idea was already clearly
described long ago in Ji Kang's aesthetic work \textit{Sheng Wu Ai Yue Lun} (On the Absence
of Sorrow and Joy in Sound)\footnote{Ji Kang [Three Kingdoms, Wei]: \textit{Sheng Wu Ai
Yue Lun}, [to be verified: edition consulted, e.g. Dai Mingyang, \textit{Collated and
Annotated Collected Works of Ji Kang}, Beijing: People's Literature Publishing House,
1962], p.~[X].}. Ji Kang holds that ``the mind and sound are manifestly two things'': sound
is an objectively existing natural acoustical phenomenon, whereas sorrow and joy are
feelings produced after the spirit has been stirred, and there is no causal relation
between the two. That is to say, music itself carries no sorrow or joy; sorrow and joy
exist in the listener.

\subsection{A Thought Experiment at the Extremes}

We take as given that the premise of music being appreciated is that the person is in a
calm state free of external interference. But human mental activity is in essence
pluralistic, and we cannot overlook the factor of extreme values. Let us perform a thought
experiment: suppose one person wins a million yuan in the lottery and another loses a
million yuan; would playing ``cheerful music'' to the one who lost money really make him
feel joy? Would playing ``sad music'' to the one who gained really make him sad? The
answer is obvious.

It should be added that the ``irrelevance'' spoken of here is more precisely expressed as
``not necessarily corresponding'': music can induce an emotional response but cannot
stipulate an emotional response---music is the occasion of emotion, not the carrier of
emotion. This tightening of the wording does not weaken Ji Kang's position; on the
contrary, it gives it firmer ground.

\subsection{A Historical Echo: The Dispute between Autonomy and Heteronomy}

In the history of Western music aesthetics this idea is no isolated case. Hanslick, in
\textit{On the Musically Beautiful}\footnote{Eduard Hanslick [Austria]: \textit{On the
Musically Beautiful: A Contribution towards the Revision of the Aesthetics of Music},
translated by Yang Yezhi, Beijing: People's Music Publishing House, [to be verified: year
of publication], p.~[X].}, published in 1854, advanced the famous theory of autonomy: the
beauty of music lies not in the expression of emotion but in ``the moving forms of tones''
themselves; against this, the emotional aesthetics of Romanticism (heteronomy) maintains
that the beauty of music lies in the expression of emotion. The two camps have argued for
nearly two hundred years without resolution. Seen through the ``stasis/motion'' layering
of this paper, the two territories in fact each find their proper place: the structural
layer is autonomous---form is beauty, stasis is quantifiable; the aesthetic layer is
heteronomous---it varies from person to person, and motion is not quantifiable. This paper
has no intention of mediating between the two positions, but only of drawing their
boundaries for them: from now on there is no need to argue further, each simply returns to
its own place.

\section{An Examination of the Quantifiability of Composition and Arrangement}

\subsection{Historical Evidence: From Zhu Zaiyu to Music Notation}

Let us return to a practical question: traditional composition and arrangement already have
so long a history---is ``composition and arrangement'' itself quantifiable? The answer can
still follow the reply given above.

From the historical angle: Zhu Zaiyu of China, in the twelfth year of the Wanli reign
(1584), divided the octave equally into twelve semitones by precise calculation, the
interval ratio of each semitone being the twelfth root of 2 (approximately 1.059463),
thereby realizing equal temperament in the true mathematical sense\footnote{Zhu Zaiyu
[Ming]: \textit{Lü Xue Xin Shuo} (A New Account of the Study of Temperament), [to be
verified: annotator, publisher, year of publication], p.~[X]; \textit{Lü Lü Jing Yi}
(Essential Principles of Pitch Pipes), [to be verified: as above or noted separately],
p.~[X].}---and this calculation was carried out by extracting roots on an eighty-one-rod
abacus, accurate to 24 decimal places. Bach's \textit{The Well-Tempered Clavier} (Volume
One, 1722)\footnote{J. S. Bach [Germany]: \textit{The Well-Tempered Clavier}, Volume One
(BWV 846--869), [to be verified: edition consulted, e.g. Leipzig: Bärenreiter, 1950],
p.~[X].} in turn confirmed the feasibility of equal temperament from the standpoint of
compositional practice. As for notation, whether it be the traditional Chinese
\textit{jianzipu} (abbreviated-character tablature) or the Western neumes, the four-line
staff, and even the five-line staff, all have labored on behalf of the quantifiability of
music. But the present author holds that we should not simply classify notation as
``for the sake of preservation and transmission''; we should return instead to quantization
and science itself---notation is in essence the symbolization of the temporal structure and
the pitch structure of music, that is, the first practical realization of quantization.

\subsection{The Mechanism of Creation: Motif, Training and Generation}

From the angle of creation: one factor truly cannot be excluded from composition---the
``motif''. In music we often attribute the ``motif'' to the composer's ``flash of
inspiration''; but is a flash of inspiration really mysticism? The present author thinks
not. This is very similar to the Transformer mechanism of modern artificial
intelligence\footnote{A. Vaswani, N. Shazeer, N. Parmar, et al., ``Attention Is All You
Need'', \emph{Advances in Neural Information Processing Systems 30 (NeurIPS 2017)}, 2017,
pp.~5998--6008.}: the human composer, too, can realize the so-called ``motif'' only after
undergoing a great deal of musical accumulation and training; even Mozart, a child prodigy
in music, could not do without a great deal of musical accumulation; and to ask someone who
has never even touched a keyboard to complete a symphony at once would almost violate the
laws of learning. From the rational angle, the so-called ``four pillars'' of
composition---harmony, musical form, counterpoint and orchestration---all have a great body
of corresponding logic, books and transmissible systems of knowledge, failing which the
whole musical fabric would fall into utter confusion; even modernist composition is a
deliberate arrangement. Indeed, twentieth-century serialism (such as Schoenberg's
twelve-tone technique and Boulez's integral serialism) pushed the quantifiability of music
to an extreme---pitch, duration, dynamics and timbre were all fully parameterized---and
the world does not on that account deny that it is music.

\subsection{The Operability of Structure and the Attribution of ``Good and Bad''}

From the mathematical angle, this can still be explained as ``the proportion of sound
occupying a given period of time''. A set of numbers is not judged to have vanished
because you replace 0 with 100; likewise, music does not become wholly invalid because you
replace a triad with a seventh chord somewhere within it. More rigorously: structural
operations do not destroy the validity of music as an organization of sound; they only
change its aesthetic effect---and the aesthetic effect belongs precisely to the
non-quantifiable layer of motion.

At this point someone is bound to ask: how, then, is one to judge whether it is good or
not? Such a question still falls within the jurisdiction of ``the aesthetics of music'' and
of ``human-made relative dynamics'', and has no connection whatever with the structure of
music itself. For instance, one cannot say that Beethoven's \textit{Fate} is necessarily
better written than \textit{The Little Painter}: structurally speaking, the musical
structure of \textit{Fate} is far more complex, but structural complexity does not equal
``good''; it depends on who the listener is---for a child, \textit{The Little Painter} obviously
has greater educational value and is more readily accepted by children. The crux of the
question returns once again to ``the aesthetics of music'' and ``human-made relative
dynamics'': because the quantified musical structure does not shift with human will, the
structure itself is neither good nor bad.

\section{The Institutional Roots of the Absence of Quantification Consciousness}

\subsection{Oral Transmission and the Dislocation of ``Right/Wrong'' Evaluation}

Since music is quantifiable, and quantification consciousness ought to be an intrinsic gene
of music, why has this question been so hard to discover for so long? Seen from the angle
of teaching methods and of social existence, music has long operated by way of ``oral
transmission and heart-to-heart instruction'', and over time the very notion of
``quantification'' has imperceptibly disappeared. Take an example: you play here, and the
teacher says ``right''; you play in another place, and the teacher says ``wrong''. But on
the argument above, musical structure itself has no right or wrong, and the aesthetic layer
is merely a thousand faces for a thousand people---so where does ``right and wrong'' come
from? It is evident that the educational system has packaged aesthetic preference as
``right and wrong'' for the purposes of assessment and grading; this is an institutional
dislocation: the structure that ought to be quantified has not been quantified, while the
aesthetics that ought not to be quantified has been quantified into a score.

\subsection{A Thought Experiment on Machine Judging}

Can the aesthetic layer, then, really be quantified? If so-called ``master'' performers
were truly taken as the reference, the best approach would seem to be to have them record
once onto a MIDI keyboard, after which all assessment would hold the student's performance
up against the machine and compute the difference. But this would plainly make artistic
expression lose its selfhood, and is utterly inadvisable. It is worth noting that this
thought experiment confirms the layering of this paper from the negative side: for a
machine to evaluate ``aesthetics'' is absurd, yet for a machine to evaluate ``structure''
is perfectly fair. The conclusion is not ``machines are unusable'' but ``machines may be
used only at the layer where they belong''.

\subsection{Disciplinary Construction and the Absence of Interdisciplinary Speculation}

From the perspective of academic and social planning: music has long been a branch of the
liberal arts and has to some extent lost the genes of the sciences, which has weakened the
capacity for interdisciplinary speculation. A clarification is needed here: this is not an
argument that music should ``become'' a science, but rather that music simultaneously
possesses a computable structural layer, and that this layer requires the participation of
scientific thinking---it is disciplinary intersection, not disciplinary replacement.
Second, when students spend a great deal of time on the adjustment of aesthetic
preferences, time being an impartial unit, their understanding at the structural layer
correspondingly decreases, and naturally there is no occasion for the possibility of
encountering additional knowledge and combining it with music.

\section{Conclusion and Prospects: Toward a Layered Quantization of Music}

\subsection{The Systematization of Four Kinds of Dislocation}

Surveying the paper as a whole, the word ``dislocation'' can run through it from beginning
to end: what this paper reveals is in fact four mutually entangled kinds of dislocation, as
shown in Table \ref{tab:cuowei}.

\begin{table}[htbp]
  \centering
  \caption{Four kinds of dislocation in the ``non-quantifiable'' work of music}
  \label{tab:cuowei}
  \begin{tabular}{p{2.6cm}p{8.0cm}p{3.2cm}}
    \toprule
    \thd{Type of dislocation} & \thd{Content} & \thd{Layer involved} \\
    \midrule
    Definitional dislocation & Conflating ``musical emotion is non-quantifiable'' with ``music is non-quantifiable'' & Propositional layer \\
    \addlinespace
    Aesthetic dislocation & Binding aesthetic value to the carrier ``human'', defining music by its producer & Ontological layer \\
    \addlinespace
    Technical dislocation & Conflating the structural layer with the aesthetic layer, either killing music with coarse quantization or abandoning quantization out of fear of difficulty & Methodological layer \\
    \addlinespace
    Institutional dislocation & Transmitting by ``oral instruction'', evaluating aesthetics by ``right/wrong'', and isolating scientific thinking by disciplinary construction & Educational and social layer \\
    \bottomrule
  \end{tabular}
\end{table}

\subsection{Prospects for a Layered Quantization Scheme}

On the basis of the argument above, the quantization of music should not be ``chopping music
to pieces'' but ``establishing a coordinate system for music''. A feasible scheme is
layered quantization: at the structural layer, every sonic event is located by absolute
temporal coordinates---the coordinates themselves lose no information, because they are
only a reference; at the aesthetic layer, the non-quantifiable breathing, deviation and
minute difference are allowed to survive explicitly in the form of ``tolerance
parameters'' between the coordinates---deviation is no longer an error but the evidence
that music is ``alive''.

The significance of this scheme for AI music arrangement lies in this: the current weakness
of music AI is not a lack of data or of computing power, but the piling of computing power
onto the wrong coordinate system. Once the structural layer obtains precise absolute
coordinates and the aesthetic layer obtains explicit tolerance expressions, AI will no
longer need to ``imitate'' music but can ``write'' it. Admittedly, this remains an
unfinished undertaking; this paper merely provides it with a foundation in ontology.

By way of conclusion, the present author would characterize himself in one sentence: this
matter is like climbing high to gaze afar---I dare not claim to have reached the summit, I
only wish that a single step's advance might spare those who come after me the trouble of
climbing. As for the edge in it, the ancients already said it before me---``I do not regret
that the ancients did not see me; I regret that the ancients did not see my wildness.''

\end{document}
